\section{胶子—胶子—规范链顶点的重正化}
最后，我们考察图~\ref{fig:div-topo-g}中拓扑 (d) 紫外发散的重正化；其图形涉及两条胶子和一条规范链，可称之为胶子—胶子—规范链顶点。与式~$(\ref{eq:gwi1})$ 的 Ward 恒等式类似，我们为“裸”场与算符构造如下 Ward 恒等式，
\begin{align}
&\langle \partial_\lambda^x A_d^\lambda(x) \partial_\rho^y A_e^\rho(y) [\Phi(\{\xi_{2z},\xi_{1z}\})]_{a b} \,  F_{b}^{\mu \nu}(\xi_{1z})  \rangle \nonumber \\
+&i \delta^{(d)}(x-y) \delta_{de} \langle [\Phi(\{\xi_{2z},\xi_{1z}\})]_{a b} \,  F_{b}^{\mu \nu}(\xi_{1z}) \rangle
\label{eq:gwi2} \\
=& g_{s} \langle f^{afg} \bar{c}_e (y) c_f (\xi_{2z}) \partial_\lambda^x A_d^\lambda(x) [\Phi(\{\xi_{2z},\xi_{1z}\})]_{g b} \,  F_{b}^{\mu \nu}(\xi_{1z}) \rangle.
\nonumber
\end{align}
式~$(\ref{eq:gwi2})$ 的图像解释见图~\ref{fig:gwi2}。与图~\ref{fig:gwi1}类似，图~\ref{fig:gwi2}左边的拓扑与图~\ref{fig:div-topo-g}中图 (d) 的相同，只是图中外胶子分别与其各自动量缩并并以坐标空间表达。在完成了前面讨论的全部重正化——包括 QCD 拉氏量、规范链和胶子—规范链顶点——之后，图~\ref{fig:gwi2}中方程的右边是 UV 有限的。也就是说，我们发现：若两条外胶子都与各自的动量缩并，则图~\ref{fig:div-topo-g}中的拓扑 (d) 没有紫外发散。

%%%%%%%%%%%%%%%%%%%%%%%%%%%%%%%
\begin{figure}[t]
	\begin{center}
		\includegraphics[width=0.55\textwidth]{images/Green-2-g.pdf}
		\caption{\label{fig:gwi2} 式~$(\ref{eq:gwi2})$ 中广义 Ward 恒等式的图像解释；虚线代表鬼场。}
	\end{center}
\end{figure}
%%%%%%%%%%%%%%%%%%%%%%%%%%%%%%%

与胶子—规范链顶点的讨论类似，Ward 恒等式帮助降低了图~\ref{fig:div-topo-g}中拓扑 (d) 的表观紫外发散。由于拓扑 (d) 的图形只有表观对数发散，来自 Ward 恒等式的这一额外压低使图~\ref{fig:div-topo-g}中的拓扑 (d) UV 有限。因此，在对拓扑 (c) 重正化之后，拓扑 (d) 不需要任何附加重正化。这与重正化 QCD 拉氏量胶子顶点的情形类似：规范不变性保证一旦三胶子顶点被重正化，四胶子顶点就没有紫外发散。
